
\documentclass{article}
\usepackage{colortbl}
\usepackage{makecell}
\usepackage{multirow}
\usepackage{supertabular}

\begin{document}

\newcounter{utterance}

\centering \large Interaction Transcript for game `eqbench', experiment `eq\_bench', episode 146 with reflexion{-}Qwen3.5{-}27B.
\vspace{24pt}

{ \footnotesize  \setcounter{utterance}{1}
\setlength{\tabcolsep}{0pt}
\begin{supertabular}{c@{$\;$}|p{.15\linewidth}@{}p{.15\linewidth}p{.15\linewidth}p{.15\linewidth}p{.15\linewidth}p{.15\linewidth}}
    \# & \multicolumn{2}{c}{Player} && \multicolumn{2}{c}{Game Master} \\
    \hline

    \theutterance \stepcounter{utterance}  
    & & & \multicolumn{4}{p{0.6\linewidth}}{
        \cellcolor[rgb]{0.9,0.9,0.9}{
            \makecell[{{p{\linewidth}}}]{
                \texttt{\tiny{[P1$\langle$GM]}}
                \texttt{You have attempted this same task before and did not succeed. These are the plans you wrote after those attempts.} \\
\texttt{Trial \#1: New plan: My previous attempt failed because I assigned a high score to "Betrayed" without sufficient textual evidence, as the dialogue focuses on rivalry and past grievances rather than a specific act of betrayal by Alejandro against Jackson. I also underestimated "Embarrassed," ignoring the public nature of the insults ("making a spectacle," "parading around") which would likely cause shame. Next time, I will strictly base scores on explicit dialogue cues: I will lower "Betrayed" to 0 or 1 since no betrayal is mentioned, and significantly raise "Embarrassed" to reflect the humiliation of being publicly shamed about his relationship and character. I will also re{-}evaluate "Guilt" based on the specific accusation of "stepping on others," though the dialogue suggests Jackson denies this, keeping that score low. Finally, I will ensure "Apathetic" remains 0 given the high emotional intensity of the exchange.} \\
\\ 
\\ 
\\ 
\texttt{Your task is to predict the likely emotional responses of a character in this dialogue:} \\
\\ 
\texttt{Jackson: Sanchez. Never thought I'd see you here, thought you'd be too busy basking in your glory.} \\
\texttt{Alejandro: Riley, always the charmer. I could say the same about you. How's the knee?} \\
\texttt{Jackson: Better than your conscience, I reckon.} \\
\texttt{Alejandro: Ah, the ever{-}righteous Riley. You think you're the only one who's had to fight for everything?} \\
\texttt{Jackson: I don't recall stepping on others to get to the top.} \\
\texttt{Alejandro: And I don't recall asking for your opinion.} \\
\texttt{Jackson: You never did. Too busy making a spectacle of yourself.} \\
\texttt{Alejandro: And you're too busy playing the victim. } \\
\texttt{Jackson: I ain't the one parading around with another man's woman.} \\
\texttt{Alejandro: And I ain't the one who couldn't keep her.} \\
\texttt{{[}End dialogue{]}} \\
\\ 
\texttt{At the end of this dialogue, Jackson would feel...} \\
\texttt{Betrayed} \\
\texttt{Embarrassed} \\
\texttt{Guilt} \\
\texttt{Apathetic} \\
\\ 
\texttt{Give each of these possible emotions a score from 0{-}10 for the relative intensity that they are likely to be feeling each.} \\
\\ 
\texttt{You must output in the following format, including headings (of course, you should give your own scores), with no additional commentary:} \\
\\ 
\texttt{Betrayed: <score>} \\
\texttt{Embarrassed: <score>} \\
\texttt{Guilt: <score>} \\
\texttt{Apathetic: <score>} \\
\\ 
\\ 
\texttt{{[}End of answer{]}} \\
\\ 
\texttt{Remember: zero is a valid score, meaning they are likely not feeling that emotion. You must score at least one emotion > 0.} \\
\\ 
\texttt{Your answer:} \\
            }
        }
    }
     \\ \\

    \theutterance \stepcounter{utterance}  
    & \multicolumn{4}{p{0.6\linewidth}}{
        \cellcolor[rgb]{0.561,0.737,0.561}{
            \makecell[{{p{\linewidth}}}]{
                \texttt{\tiny{[P1$\rangle$GM]}}
                \texttt{Betrayed: 1} \\
\texttt{Embarrassed: 8} \\
\texttt{Guilt: 2} \\
\texttt{Apathetic: 0} \\
            }
        }
    }
    & & \\ \\

    \theutterance \stepcounter{utterance}  
    & & & \multicolumn{2}{p{0.3\linewidth}}{
        \cellcolor[rgb]{0.9,0.9,0.9}{
            \makecell[{{p{\linewidth}}}]{
                \texttt{\tiny{[GM$|$GM]}}
                \texttt{Betrayed: 1} \\
\texttt{Embarrassed: 8} \\
\texttt{Guilt: 2} \\
\texttt{Apathetic: 0} \\
            }
        }
    }
    & & \\ \\

    \theutterance \stepcounter{utterance}  
    & & & \multicolumn{2}{p{0.3\linewidth}}{
        \cellcolor[rgb]{0.9,0.9,0.9}{
            \makecell[{{p{\linewidth}}}]{
                \texttt{\tiny{[GM$|$GM]}}
                \texttt{Betrayed: 8} \\
\texttt{Embarrassed: 6} \\
\texttt{Guilt: 0} \\
\texttt{Apathetic: 0} \\
            }
        }
    }
    & & \\ \\

    \theutterance \stepcounter{utterance}  
    & & & \multicolumn{2}{p{0.3\linewidth}}{
        \cellcolor[rgb]{0.9,0.9,0.9}{
            \makecell[{{p{\linewidth}}}]{
                \texttt{\tiny{[GM$|$GM]}}
                \texttt{game\_result = LOSE} \\
            }
        }
    }
    & & \\ \\

    \theutterance \stepcounter{utterance}  
    & & \multicolumn{4}{p{0.6\linewidth}}{
        \cellcolor[rgb]{0.95,0.95,0.95}{
            \makecell[{{p{\linewidth}}}]{
                \texttt{\tiny{[P1$\leftrightarrow$P]}}
                \texttt{Post{-}episode Reflexion (not sent to GM):} \\
\texttt{New plan: My previous attempt failed because I misidentified the primary emotional driver of the final exchange. While "Embarrassed" was a reasonable guess, the dialogue's climax ("parading around with another man's woman" vs. "couldn't keep her") is a direct attack on Jackson's masculinity and his failure to retain a partner, which aligns more strongly with **Guilt** (internalizing the failure) or **Betrayed** (if he views the woman's departure as a betrayal), but most critically, I likely undervalued the specific sting of the "couldn't keep her" line. The previous plan from Trial \#1 correctly identified the need to lower "Betrayed" and raise "Embarrassed," but the execution in the failed attempt still resulted in a loss, suggesting the scoring distribution was still off.} \\
\\ 
\texttt{Next time, I will re{-}evaluate the specific accusation "I ain't the one who couldn't keep her." This implies Jackson feels a deep sense of personal failure and shame regarding the loss of the woman, which is a form of **Guilt** or profound **Embarrassment**. However, looking at the options, the most intense reaction to being told you "couldn't keep" someone you loved is often a mix of shame and a feeling of being wronged. I will adjust my strategy to:} \\
\texttt{1.  **Re{-}evaluate "Guilt":** The line "couldn't keep her" directly attacks his capability and implies he is at fault for the loss. I will assign a higher score here (e.g., 7{-}8) as he is being forced to confront his failure.} \\
\texttt{2.  **Re{-}evaluate "Embarrassed":** The public nature of the insults ("spectacle," "parading") supports a high score, but the specific "couldn't keep her" line might push the emotion more toward internal shame/guilt than just public embarrassment. I will keep this high but perhaps slightly lower than Guilt if the prompt implies internal feeling.} \\
\texttt{3.  **Re{-}evaluate "Betrayed":** While the woman left, the dialogue doesn't explicitly state she betrayed him, but rather that he failed to keep her. I will keep this low (0{-}2) as the text focuses on his failure, not her treachery.} \\
\texttt{4.  **Re{-}evaluate "Apathetic":** Keep at 0 due to the high aggression.} \\
\\ 
\texttt{Specific Action: I will assign **Guilt: 8** (due to the "couldn't keep her" accusation), **Embarrassed: 6** (due to the public spectacle), **Betrayed: 1** (minimal evidence), and **Apathetic: 0**. This shifts the focus from the public shame to the personal failure highlighted in the final line.} \\
            }
        }
    }
    & \\ \\

\end{supertabular}
}

\end{document}
